Motivate tariffication and the selection of relevant criteria

Insurer: Serves collective of policyholders
Has to cover total claims
Have discussed: (random) total claim size \( S \mapsto \pi(S) \)
Individual or collective model of risk explain \( S \)
Question: How do we split \( \pi(S) \) among policyholders?
Simple answer: Homogeneous insurance portfolio → individual premium contribution
\( \frac{\pi(S)}{\text{# insured risks}} \)
Problem: Homogeneity of collective fails
Insurance sums differ, Individual behaviour of policyholders more or less risk seeking

→ Tariffication: individual premium contribution should reflect quality of specific
insurance policy (e.g. individual risk profile)

Goal: Sort insurance contracts in a number of different groups — risk cells or risk classes
Within each group: Policyholders share similar traits/features
Groups not too small! (“Law of large numbers should apply”)
→ (approximately) homogeneous individual claim size distributions
Individual premium contribution uniform in each risk cell

Important distinction:
tariff criteria vs. risk criteria
Risk criteria: Features of policyholder/insured object with predictive value for risk a
specific contract poses/claim size it produces
Tariff criteria: Features of policyholder/insured object actually used for tariffication
Ideally: Tariff criteria selected from risk criteria (but this isn't always the case)

Statistical significance?
Are risk criteria uncorrelated or correlated, therefore redundant?
Is insurer allowed to use risk criterion as discriminatory tool in tariffication?


Formulate the framework of tariffication quantitatively
and describe the method of marginal averages.


For \(\ell \in \mathbb{N}: [\ell] := \{1, 2, ..., \ell\}\)
\(r \in \mathbb{N} \widehat{=}\) number of tariff criteria
\(M_k \widehat{=}\) kth tariff criterion, \(k \in [r]\)
\(n_k \in \mathbb{N} \widehat{=}\) number of characteristics of \(M_k\)
\(a_{jk} \widehat{=}\) jth characteristic of \(M_k,\ j \in [n_k]\)

Vector \((j_1, ..., j_r) \in \prod_{k=1}^r [n_k] \widehat{=}\) particular risk cell
\(t := \prod_{k=1}^r n_k \widehat{=}\) total number of risk cells
\(n := \sum_{k=1}^r n_k \widehat{=}\) total number of characteristics of tariff criteria
\(\bar{b} \widehat{=}\) basic premium
\(b_{j_1, ..., j_r} \widehat{=}\) individual premium contribution in risk cell \((j_1, ..., j_r)\)

In a mult. model the marginal factors \( u_{kj_k}_{k \i [r],j_k \in [n_k]} \) s.t.
\( b_{j_1, \dots, j_r} = \bar{b}\cdot\prod^r_{k=1} u_{kj_k} \)

Shall focus on the case of two tariff criteria, \(M_1\) and \(M_2\) (\(r=2\))
\(n_1 := I \in \mathbb{N}\), \(n_2 := K \in \mathbb{N}\)

Computation of marginal factors uses data about claims and volume
Fix cell \((i, j) \in [I] \times [K]\)
\(s_{ij} > 0\) total claims in cell \((i, j)\)
\(v_{ij} > 0\) total volume in cell \((i, j)\)
\(v_{ij} \widehat{=}\) total insurance sum or units per year

Abbreviate
\(
x_i := u_{1i}, \quad y_j := u_{2j}, \quad (i, j) \in [I] \times [K].
\)

Shall consider volume-adjusted quantity
\[
z_{ij} := \frac{s_{ij}}{v_{ij}}, \quad (i, j) \in [I] \times [K],
\]
called claims requirement (in case of units per year) or claims rate (in case of insurance sums), \((i, j) \in [I] \times [K]\)
Set
\[
\mu := \frac{\sum_{(i, j) \in [I] \times [K]} s_{ij}}{\sum_{(i, j) \in [I] \times [K]} v_{ij}} \widehat{=} \frac{\text{total losses}}{\text{total volume}}.
\]
This shall serve as basic premium: \(\bar{b} = \mu\)
Akin to pure risk premium (expected claim size)

The method of marginal averages prescribes the following marginal factors:
\[
x_i := \frac{\sum_{j=1}^K s_{ij}}{\mu \sum_{j=1}^K v_{ij}}, \quad i \in [I],
\]

\[
y_j := \frac{\sum_{i=1}^I s_{ij}}{\mu \sum_{i=1}^I v_{ij}}, \quad j \in [J].
\]

The resulting individual premium charged in cell \((i, j) \in [I] \times [J]\) is given by
\[
b_{ij} = \mu x_i y_j.
\]


State the method of total marginal sums from Bailey & Jung

The method of marginal averages has a tendency for over/underestimation,
a potential solution is given by the alternative marginal factors:

\( x_i = \frac{\sum^K_{j=1} s_{ij}}{\mu \sum_{j=1}^K v_{ij}y_j}, i \in [I] \)
\( y_j = \frac{\sum^l_{i=1} s_{ij}}{\mu \sum_{i=1}^l v_{ij}x_j}, j \in [K] \)

Usually numerical methods are required to solve this system of equations,
by hand this can be mimicked by choosing initial \( y_j \) (hopefully close to the actual solution)
determining \( x_i \), with them new \( y_j \) and repeating this iteratively until the value stabilizes.


Build up the framework for Bonus-Malus Systems (BMS) via Markov Chains

Primary differentiation of premia: 
Using risk cells, collective is sorted into
approximately homogeneous subcollectives
Secondary differentiation of premia: 
Adjustment of premia according to individual loss experiences

BMS corresponds to the secondary case.
Advantage: Incentivise risk-averse behaviour of policyholders
Reward absence of losses in one period with lower premium in subsequent one,
Penalise losses in one period with higher premium in subsequent one.

Let there be \( m \)-BM cells (or classes), \( \mathcal{C} = \{c_0, \dots, c_{m-1}\} \)
with \( c_1 < \dots < c_m \) and let \( \{X_t\}_{t \in \mathbb{N}} \) be the classification process
for an individual.
Initially put all individuals into one cell e.g. \( X_0 = c_0 \).
The BM rule assigns new cells \( \phi : \mathcal{C} \times \mathbb{N}_0 \to \mathcal{C} \)
\( \forall k, k' \in \mathbb{N}_0 : k \leq k' \implies \phi(c, k') \leq \phi(c, k) \),
the corresponding premium is \( \mu \pi(c_i) \), where \( \mu \) is some base premium.

A (discrete) Markov chain is characterized by two objects
The matrix of transition probabilities 
\( Q(t) = Q^t : \mathbb{R}^{|\mathcal{C}|} \times \mathbb{R}^{|\mathcal{C}|}\)
with entries
\( q_{cc'}(t)_{(c,c') \in \mathcal{C} \times \mathcal{C}} = \mathbb{P}[X_{t+1} = c' |X_t = c] \),
if \( \mathbb{P}[X_t = c] > 0 \) and otherwise \( 0\).
The initial distribution \( p := \mathbb{P}_{X_0} \)

It holds
\( \sum_{j \in \mathcal{C}} q_{ij} = 1 \) for all \( i \).
\( \mathbb{P}[X_t = c] = (p^TQ^t)_c = \sum_{i=1}^m p(i)(Q^t)_{ic} \)

Let \( N_t \) be the number of claims with individual claim sizes \( Y_i^t \)
then define \( Q^{(k)}_{c,c'} = \mathbf{1}_{\phi(c,k)}(c') \), it holds
that \( Q = \sum_{k=0}^\infty \mathbb{P}[N_0 = k] Q^{(k)} \) defines a transition matrix.

A typical example is \( \phi(c, 0) = c + 1 \wedge c_m \) and \( \phi(c, k) = c - 1 \vee 1 \) for \( k > 0 \),
e.g.


State methods to determine the base premium for BMS

Let \( Z_t = \mu\pi(X_t) \), then 
\( \mathbb{E}[Z_t] = \mu \sum{i=1}^m \pi(i)\mathbb{P}[X_t = i] = \mu \sum_{(c,i)}^m \pi(i)p(c)Q^t_{c,i} \)
The approaches are
1. \( \mathbb{E}[Z_u] = \mathbb{E}[S_u] \) for fixed \( u \).
2. \( \sum_{t=0}^{n-1} \mathbb{E}[Z_t] = \sum_{t=0}^{n-1}\mathbb{E}[S_t] \)
in the second \( \mu \) becomes dependend on the number of observations \( n \).

Let \( \alpha = \mathbb{E}[N_u], \beta = \mathbb{E}[Y^{0}_1] \)

Consider the basic premia \(\mu_u\), \(u \in \mathbb{N}_0\), defined by approach 1. and \(\mu(n)\), \(n \in \mathbb{N}\), from approach 2.
If \(p\) is a \(Q\)-invariant initial distribution, i.e. \(p^T = p^T Q\), then
\(
\forall n \in \mathbb{N} \forall u \in \mathbb{N}_0 : \mu_u = \mu(n) = \mu(1) = \frac{\alpha \beta}{\sum_{c=1}^m p(c) \pi(c)}.
\)
If \(Q\) is weakly ergodic, i.e. there is \(p_\infty \in \mathbb{R}^m\) such that, for all unit vectors \(e_c\), \(c \in [m]\),
\(
\lim_{t \to \infty} e_c^T Q^t = p_\infty^T,
\)
we have
\[
\lim_{u \to \infty} \mu_u = \lim_{n \to \infty} \mu(n) = \frac{\alpha \beta}{\sum_{c=1}^m p_\infty(c) \pi(c)}.
\]


Describe the relevant steps for constructing a mortality table. 
Assuming stationarity, explain how survival probabilities can be derived.
Name an approach for modelling the remaining lifespan analytically.

Derivation of mortality tables:
• The population is first divided according to certain risk characteristics: Age,
gender, region, smoker or non-smoker, . . . (The more risk characteristics are
taken into account, the smaller the homogeneous subclasses (sample size).)
• Mortality tables are derived in a three-stage process:
– Estimate the raw mortality probabilities
– smoothing of raw mortality probabilities
– Estimation of safety margins
• In practice, mortality trends must also be recognised and predicted accurately.
Condition for a stationary process: In practice, it is often assumed that the condition
for a stationary process
\( \mathbb{P}[T_{x+s} > t] = \mathbb{P}[Tx > s + t| Tx > s] \), \( s, t, x \geq 0 \)
is fulfilled for the distribution of future lifespans. 
Implication: In particular, the condition for a stationary process implies
\( {}_t p_x = \mathbb{P}[T_x > t] = \mathbb{P}[T_0 > x + t|T_0 > x] = \frac{{}_{t+x}p_0}{{}_{x} p_0} \)

which means that the distribution of future lifespans of an \( x \)-year old person is
already fixed by the distribution of future lifespans of a newborn.

Assumption: The future lifetime follows a natural law (mortality law) that can be
expressed through a simple formula for the distribution function \( G_x \) with only a few
parameters.
De Moivre:
\( \mu(x+t) = (\omega - x - t)^{-1}, g_x(t) = (\omega - x)^{-1}, G_x(t) = tg_x(t) \), \( 0 \leq x < \omega \)
Has a maximum age \( \omega \) and distributes \( T_x \) uniformly on \( [0, \omega-x] \) and mortality increases with \( t \).
Weibull
\( \mu(x+t) = k(x+t)^n, s(x) = e^{-ux^{n+1}} \), \( k > 0, n > 0, x \geq 0 \)

Define the total losses of an insurance company using an individual model of risk.
What assumptions does a homogeneous portfolio satisfy? How does the collective
model differ from the individual model?

total claim size: \( S = \sum_{i=1}^n Y_i \) with random individual claims \( Y_1, \dots, Y_n \) and
deterministic number of risks n in the portfolio (Typically, \( \mathbb{P}[Y_i = 0] > 0 \)P
(no claim with pos. probability))
• homogeneous portfolio: individual claim sizes \( Y_i \) are independent and
identically distributed
• Distinctive features of collective model: Only consider risks that produce claims
and acknowledge their number is random

Nennen Sie den Zweck von Bonus-Malus-Systemen bei der Gestaltung von Prämien.

Bonus-Malus-Systeme dienen der Prämiendifferenzierung basierend auf beobachteter Schadenvergangenheit, 
d.h. gute / schadenfreie Risiken zahlen eine günstigere Prämie als schlechtere / schadenbehaftete Risiken.



What are assumptions for fractional ages or interpolations of life table entries?


Let \( t \in [0, 1] \) and \( x \in \mathbb{N}_0 \), to interpolate the survival probabilities of discrete entries of a life
table the following  assumptions are common:  
1. Linear interpolation: 
\( S(x + t) = (1 - t) S(x) + t S(x + 1) \). 
This is known as the uniform distribution or, perhaps more properly, a uniform distribution of deaths assumption within each year of age. Under this assumption \( {}_t p_x \) is a linear function.  
2. Exponential interpolation, or linear interpolation on 
\( \log S(x + t) \): \(\log S(x + t) = (1 - t) \log S(x) + t \log S(x + 1).\) 
This is consistent with the assumption of a constant force of mortality within each year of age. Under this assumption \( {}_t p_x \) is exponential.  
3. Harmonic interpolation: 
\(\frac{1}{S(x + t)} = (1 - t) \frac{1}{S(x)} + t \frac{1}{S(x + 1)}.\) 
This is what is known as the hyperbolic (historically \textit{Balducci}) assumption, for under it \( \, ^t p_x \) is a hyperbolic curve.

One can then derive the remaining magnitudes e.g. for the uniform assumption we have

\( {}_t q_x = \frac{S(x) - S(x+t}{S(x)} = \frac{S(x) - ((1-t)S(x) + tS(x+1))}{S(x)} = \frac{tS(x) - tS(x+1)}{S(x)} = tq_x \)

The corresponding forces of mortality are given by
1. \( \mu(x+t) = \frac{q_x}{1 - tq_x} \)
2. \( \mu(x+t) = - \log p_x \)
3. \( \mu(x+t) = \frac{q_x}{1 - (1-t)q_x} \)


